% =============================================================
% RAR — Draft 0 Introduction (DISPOSABLE SCAFFOLD)
% Purpose: set guardrails for the evaluation. NOT the final intro.
% The final introduction is rewritten from scratch AFTER the evaluation
% (Principle 1, introduction-twice). Draft 0 is reference material only.
% Mechanism: SOFT routing via reconstruction-temperature-weighted dual heads.
% Last updated: 2026-07-09
% =============================================================

\section{Introduction}

% M1 Stakes: who suffers, why the domain matters (problem-first, Principle 7)
Encrypted-traffic classifiers identify applications and threats over traffic that no longer exposes payloads, and the set of classes they must recognize grows continuously.
New applications launch, new malware families appear, and each addition forces the classifier to recognize a class it has never seen.
Retraining the classifier from scratch over the full class history on every addition is expensive, and the cost grows with each new class.
Operators need a classifier that absorbs new application and threat classes while retaining the classes it already knows.

% M2 Problem Gap: three structural limitations of prior CIL methods (named baselines)
Class-incremental learning (CIL) targets this need by adding classes without retraining on the full history.
Existing CIL methods, including LwF, iCaRL, BiC, and DER, mitigate forgetting through two mechanisms: replaying a small buffer of historical samples and distilling the old model into the new one.
Applied to encrypted traffic, these methods leave old classes forgotten, for three structural reasons.
First, expanding the classification head with new-class outputs makes old and new classes compete for one shared decision boundary.
Second, uniform distillation treats every sample identically, regardless of whether it carries old or new knowledge.
Third, replay depends on a limited buffer that cannot represent the full old-class distribution.

% M3 Key Abstraction: reconstruction error -> sample-level knowledge ownership (observation, not inevitability)
A module reconstructed from the old model and its replay buffer tends to reproduce old-class samples with lower reconstruction error and new-class samples with higher reconstruction error.
Reconstruction error therefore provides a sample-level knowledge ownership signal: for each input, the probability that the input belongs to old versus new knowledge.
We call this per-input signal \emph{knowledge ownership estimation}.

% M4 Design Intuition: Knowledge Ownership Estimation -> Soft Routing -> Dual Knowledge Paths -> Calibration
RAR (Reconstruction-Aware Routing) turns the knowledge ownership signal into a soft routing weight over two knowledge paths, an old-knowledge head and a new-knowledge head.
Both heads contribute to every prediction, weighted by the sample's estimated old/new knowledge probability.
A reconstruction temperature $\tau$ controls routing softness: small $\tau$ approaches deterministic head selection, large $\tau$ smoothly fuses the two heads.
RAR then calibrates the combined logits into a joint prediction over all classes.
The probability separates old from new knowledge softly, and calibration recombines the two paths.
The design follows a single line: knowledge ownership estimation, soft routing, dual knowledge paths, calibration.

% M5 Contributions: 4 numbered, claim-first
We make four contributions.
\begin{enumerate}
  \item \textbf{Sample-level knowledge ownership estimation.} We estimate each sample's old/new knowledge probability from reconstruction error, enabling sample-aware knowledge balancing beyond sample-agnostic knowledge transfer.
  \item \textbf{Dual-head knowledge path decoupling.} We route old and new knowledge paths to separate heads, removing the shared-classifier competition that causes forgetting.
  \item \textbf{Logit calibration over dual knowledge paths.} We calibrate the two heads' outputs into a unified classification space, raising mixed-class accuracy over an uncalibrated dual-head baseline.
  \item \textbf{Validation across six encrypted-traffic datasets.} We show RAR reduces forgetting and improves mixed accuracy over ET-BERT-based incremental learning baselines and representative replay-based CIL methods on six public datasets (NUDT-Mobile, CIC-DoHBrw, CIC-IoT2023, CSTNET-TLS1.3, ISCX-Tor, EBSNN-WebApp).
\end{enumerate}

% M6 Results Preview: BLANK in Draft 0. Fill after evaluation.
% Tentative headline: RAR cuts forgetting 53-67% and raises mixed-class accuracy up to 11.3% over ET-BERT across six datasets.
